\documentclass[11pt]{article}
\usepackage{mathblatt}
% Lösungen nach mathe-nachhilfe gemeinsam.md „Lösungen“ und pruefung.md § 6.
\newcommand{\kennung}{T2U}
\begin{document}
\pfheftstil
\pfheftkopf[\kennung]{Winkel ohne Rechnung bestimmen oder begründen}{P10 \textperiodcentered{} Lösungen}

\begin{pfloesung}{}
\lz{1.}{B}{Raute ist ein Parallelogramm: gegenüberliegende Winkel gleich; benachbarte ergeben zusammen $180^\circ$}{}
\end{pfloesung}
\begin{pfloesung}{}
\lz{2.}{B}{Trapez: mindestens ein Paar paralleler Seiten}{}
\end{pfloesung}
\begin{pfloesung}{}
\lz{3a)}{$57^\circ$}{$180^\circ - 52^\circ - 71^\circ$}{}
\lz{b)}{$76^\circ$}{$360^\circ - 84^\circ - 97^\circ - 103^\circ$}{}
\end{pfloesung}
\begin{pfloesung}{}
\lz{4.}{$\varepsilon = 38^\circ$}{Dreieck $ADC$, rechter Winkel bei $D$: $180^\circ - 90^\circ - 52^\circ$}{}
\end{pfloesung}
\begin{pfloesung}{}
\lz{5.}{$\beta = 105^\circ$}{benachbarte Winkel im Parallelogramm: $180^\circ - 75^\circ$}{}
\end{pfloesung}
\begin{pfloesung}{}
\lz{6.}{$\gamma_1 \approx 58{,}3^\circ$}{Dreieck $ADC$, rechter Winkel bei $D$: $180^\circ - 90^\circ - 31{,}7^\circ$}{}
\end{pfloesung}
\begin{pfloesung}{}
\lz{7.}{Lea hat recht: $\delta = 55^\circ$}{Dreieck $ABC$: $\angle ACB = 180^\circ - 90^\circ - 30^\circ = 60^\circ$; Nebenwinkel $\angle ACD = 120^\circ$; $\delta = 180^\circ - 5^\circ - 120^\circ = 55^\circ$}{}
\end{pfloesung}
\begin{pfloesung}{}
\lz{8.}{$\overline{BC} = 5$ cm}{gleiche Winkel bei $A$ und $B$ $\Rightarrow$ gegenüberliegende Seiten gleich lang}{}
\end{pfloesung}
\begin{pfloesung}{}
\lz{9.}{$\overline{BC} = 6$ cm; $\gamma = 40^\circ$}{gleiche Winkel bei $A$ und $B$ $\Rightarrow$ $\overline{BC} = \overline{AC}$; $\gamma = 180^\circ - 2 \cdot 70^\circ$}{}
\end{pfloesung}
\begin{pfloesung}{}
\lz{10.}{gleichschenklig}{Winkel bei $A$ und bei $B$ im Dreieck $ABF$ je $45^\circ$ (bei $B$ auch: $180^\circ - 90^\circ - 45^\circ$) $\Rightarrow$ $\overline{AF} = \overline{BF}$}{}
\end{pfloesung}
\begin{pfloesung}{}
\lz{11.}{rechter Winkel}{Dreiecke $ADC$ und $DBC$ gleichschenklig mit $90^\circ$ bei $D$: Basiswinkel je $45^\circ$; bei $C$: $45^\circ + 45^\circ = 90^\circ$}{}
\end{pfloesung}
\begin{pfloesung}{}
\lz{12.}{$\gamma = 90^\circ$}{$\overline{BD}$ halbiert $\overline{AC}$: $\overline{AF} = \overline{FC} = 25$ cm $= \overline{DF}$; Dreiecke $AFD$ und $CFD$ gleichschenklig-rechtwinklig: je $45^\circ$ bei $D$; $45^\circ + 45^\circ = 90^\circ$}{}
\end{pfloesung}
\end{document}
